\section{Roadmap: the milestones ahead}
\label{sec:roadmap}

This document is a living artifact. \textbf{Milestones 1 through 5 are complete} (pending PI sign-off) ---
the measurement map, its stratification, their temporal coherence, their prognostic value, and the treatment
boundary --- each a separate milestone with its own acceptance gate, none begun before the previous was
locked. \textbf{PI sign-off is next; a true treatment-selection study (M5b) needs randomized data.}

\milestonestrip{5}

\paragraph{M2 --- Validated strata.} \statuspill{facekr}{COMPLETE}\ \emph{(Section~\ref{sec:strata}).} On the
nine-dimension coordinates the transdiagnostic space is a \textbf{continuum, not biotypes}: eight extreme
phenotypes (every patient a soft blend), severity as the spine, distinct metabolic and inflammatory corners,
and a partition that cuts across DSM-5 (adjusted Rand index $\approx 0$) while describing the coordinates far
better than diagnosis ($\approx 21\%$ vs $\approx 5\%$ variance explained). Internal/descriptive validity only
--- the predictive and treatment head-to-heads are M4/M5.

\paragraph{M3 --- Temporal coherence.} \statuspill{facekr}{COMPLETE}\ \emph{(Section~\ref{sec:temporal}).} The
fixed map and strata are scored forward onto follow-up (V0\,$\to$\,V1\,$\to$\,V2; fit-once, score-anyone), and
the signal is shown to \emph{persist}. This is where invariant~I3 pays off: because the map was fixed on
baseline, later visits are an honest out-of-sample test. The measurement is invariant across visits (five of
six backbone axes), and the M2 geometry replays --- the cohort slides along the severity spine (state) while
individual biological positions and archetype identity are preserved (trait), confirmed by a variance
decomposition and an independent geometric test that agree. The operative read-out, carried into M4:
\emph{stratify on the durable biology (cognition, metabolic, inflammatory); monitor the moving symptoms
(severity, suicidality, sleep)}. Persistence is not prediction --- that was M4.

\paragraph{M4 --- Prognosis.} \statuspill{facekr}{COMPLETE}\ \emph{(Section~\ref{sec:prognosis}).} The
biology$\perp$severity bet paid off, for \emph{functioning}: on the fixed objects, an errors-in-variables
nested head-to-head (diagnosis $+$ severity $+$ baseline outcome vs.\ $+$ the map) shows the durable
metabolic/inflammatory axes and the eight archetypes predict two-year functioning incrementally
($\Delta$ELPD $+46$; functional remission $14\%$--$60\%$ across archetypes), robust to attrition, reliability
and a permutation null. The signal is \emph{co-informative} with the DSM-5 subtypes (complementary) and
\emph{course-dependent} (episodic bipolar/depression, not baseline-saturated schizophrenia); severity itself
is autoregression-determined. The map's worth is group-level stratification plus continuous functional
forecasting, not a large individual-binary boost.

\paragraph{M5 --- Treatment.} \statuspill{facekr}{COMPLETE}\ \emph{(Section~\ref{sec:treatment}).} The
data-feasibility check overturned the initial reading: treatment data exists in the per-cohort source
dictionaries (unharmonized), and was mapped to common drug-class exposures. A proper causal pipeline
(overlap gate $\to$ propensity $\to$ doubly-robust errors-in-variables moderation $+$ E-value) shows that,
on observational treatment-as-usual, map position does \emph{not} reliably \emph{moderate} response: lithium
in bipolar (the cleanest contrast) is a well-powered \emph{null}, an antipsychotic metabolic/inflammatory
signal is suggestive but unconfirmed, and clozapine is confounded by channeling. The prognosis
\emph{survives} treatment adjustment (strengthening M4). The boundary is \emph{earned}, not assumed ---
treatment \emph{selection} needs the randomized/trial-arm data of a future \textbf{M5b}.

\subsection{Milestones 1--5 are complete --- and the small follow-ons they leave}
The deliverables are all in hand: the certified nine-dimension map with its hardening chain
(Sections~\ref{sec:validation}--\ref{sec:scoring}) and prior$\to$posterior atlas; on those coordinates, the
validated transdiagnostic continuum --- eight extreme phenotypes and a soft tessellation, internally validated
and adjudicated against DSM-5 (Section~\ref{sec:strata}); and the demonstration that both \emph{persist} across
follow-up --- invariant measurement and a replaying spine-versus-corner geometry (Section~\ref{sec:temporal});
and the demonstration that the durable axes \emph{predict} two-year functioning incrementally beyond diagnosis
and severity (Section~\ref{sec:prognosis}); and, on observational treatment-as-usual, the \emph{earned}
finding that the map does not reliably moderate treatment, while the prognosis survives treatment adjustment
(Section~\ref{sec:treatment}). What remains is a short, principled list of
follow-ons, none of which changes a conclusion: full-$N$ full-precision certification of the explicit-latent
block (the suicidality$\sim$developmental $\Phi$ coupling); full-$N$ projection of the non-Gaussian per-patient
scores; a separate bootstrap and correlated-$G$ arm for mania/substance (they already carry the 9-dim
cross-seed $\varphi$ $0.993$); and a hurdle likelihood for \code{isf09a} if its count precision is needed.
\textbf{PI sign-off on the adjudication, the atlas, and the strata/temporal/prognosis/treatment findings
locks M1\,$+$\,M2\,$+$\,M3\,$+$\,M4\,$+$\,M5; a true treatment-selection study (M5b) needs randomized data.}
